\subsection{对线性系统的应用}

命题 6. (i) 给定一个齐次线性方程组

\[
\sum_ {\iota \in I} \alpha_ {\mu \iota} \xi_ {\iota} = 0 (\mu \in M)\tag{2}
\]

其系数 \(\alpha_{\mu \iota}\) 属于 \(K'\) ，该方程组的每个由 K 中元素组成的解（ \(\xi_{\iota}\) ）都是 (2) 的解（ \(\xi_{\iota}'\) ，由 \(K'\) 中元素组成的）的、以 K 中系数作出的线性组合。

\begin{enumerate}
\def\labelenumi{(\roman{enumi})}
\setcounter{enumi}{1}
\tightlist
\item
  给定一个线性方程组
\end{enumerate}

\[
\sum_ {\iota \in I} \alpha_ {\mu \iota} \xi_ {\iota} = \beta_ {\mu} \quad (\mu \in \mathbf {M})\tag{3}
\]

其系数 \(\alpha_{\mu \iota}\) 和右端项 \(\beta_{\mu}\) 属于 \(K'\) ，如果存在该方程组的一个由 K 中元素组成的解，则也存在一个由 \(K'\) 中元素组成的解。

\begin{enumerate}
\def\labelenumi{(\roman{enumi})}
\item
  对每个集合 S，设右向量 K-空间 \(\mathrm{K}_{d}^{(\mathrm{S})}\) 被赋予 \(K^{\prime}\) -结构 \(\mathrm{K}_{d}^{\prime(\mathrm{S})}\) 。设 f 是把 \(\mathrm{K}_{d}^{(1)}\) 映入 \(\mathrm{K}_{d}^{(\mathrm{M})}\) 的 K-线性映射，它把每个向量 \((\xi_{\iota})_{\iota\in\mathbb{I}}\) 映为向量 \((\zeta_{\mu})_{\mu\in\mathbb{M}}\) ，其中 \(\zeta_{\mu}=\sum_{\iota\in\mathbb{I}}\alpha_{\mu\iota}\xi_{\iota}\) 对所有 \(\mu\in M\) 成立。显然 f 在 \(K^{\prime}\) 上有理；它的核 V（即 K 中方程组 (2) 的解集）是 \(\mathrm{K}_{d}^{(\mathrm{I})}\) 的一个在 \(K^{\prime}\) 上有理的子空间（第 3 号，命题 3 的推论 2），因而由 (2) 的在 \(K^{\prime}\) 中的解生成。
\item
  我们把 K 视为一个左向量 \(\mathbf{K}'\) -空间；存在一个 \(\mathbf{K}'\) -线性投影算子 \(p\) ，它把 K 映到其向量子空间 \(\mathbf{K}_s'\) 上（§ 7，第 3 号，命题 4）；若 \((\xi_{\iota})\) 是 (3) 在 K 中的一个解，则 \(\sum_{\iota \in I} \alpha_{\mu \iota} p(\xi_{\iota}) = p\left(\sum_{\iota \in I} \alpha_{\mu \iota} \xi_{\iota}\right) = p(\beta_\mu) = \beta_\mu\) ，这就证明了 \((p(\xi_{\iota}))\) 是 (3) 在 \(\mathbf{K}'\) 中的一个解。
\end{enumerate}

一个环 K 称为在子环 \(K'\) 上（左）忠实平坦的，如果命题 6 对 K 和 \(K'\) 成立；我们将在后面详细研究这一概念（交换代数，I，§3）。

